% GENERATED FILE. Edit canonical lessons, then run npm run generate:books.
% canonical-source-hash: fnv1a64:d882c1e89d4c56b0
% canonical-lessons: ES-C369-rio, ES-C369-arena, ES-C369-polvo, ES-C369-humo, ES-C369-ceniza

\chapter{The River, The Sand And The Ash}
\label{ch:el-rio-y-la-arena}
\hlchaptermodality{369}

\begin{chapteropening}
\textbf{By the end of this chapter:} \emph{I can say el río, la arena, el polvo, el humo and la ceniza, and name what the ground carries and what a fire leaves.}

Complete ES-C369-ceniza: name what the ground carries and what a fire leaves behind.
\end{chapteropening}

\section[el río]{el río — and the English river that is not the same word}

\label{lesson:ES-C369-rio}



\begin{quote}
Say \emph{the wolf}, then \emph{the bull}. (\emph{El lobo}; \emph{el toro}.)
\end{quote}

\subsection*{What to know first}
\begin{itemize}
\raggedright
  \item el agua — what runs in it.
  \item el valle — the shape it cuts.
\end{itemize}

\begin{sounds}
\begin{itemize}
\raggedright
  \item \emph{RÍ-o}, two beats with the weight on the first. The \emph{r} at the start of a word is rolled, and the written accent keeps the \emph{i} and the \emph{o} apart as two separate beats. $\to$ pronunciation reference
\end{itemize}
\end{sounds}

\begin{cousinweb}
\textbf{el río} — \textquotedblleft{}the river.\textquotedblright{}

\emph{Río} is Latin \textbf{rīvus}, a stream. English has two words from it, and both still carry the water:

\begin{itemize}
\raggedright
  \item a \textbf{rival} is \emph{rīvālis}, one of the people drawing on the \textbf{same stream}. Share a watercourse with a neighbour and you will eventually disagree about it, and the Romans built that observation into a noun;
  \item to \textbf{derive} is \emph{dērīvāre}, to lead water \textbf{off} a channel. A derived word has been drawn off from an earlier one the way a ditch is drawn off a river.
\end{itemize}

Now the part worth slowing down for. English \textbf{river} is \textbf{not} this word.

\emph{River} came through Old French from a different Latin noun, \textbf{rīpa}, which means a \textbf{bank} — the edge, not the water. The French word named the land along the water first and slid across to the water itself later. And that same \emph{rīpa} sits inside \textbf{arrive}, which began as \emph{ad rīpam}, \textquotedblleft{}to the shore\textquotedblright{}: to arrive was to reach the bank.

So English \textbf{river} belongs with \textbf{arrive}, and Spanish \textbf{río} belongs with \textbf{rival}. Two words that look alike, mean almost the same thing, and come from two Latin nouns with no relationship to each other at all.

This is the trap the whole business of etymology exists to catch. Resemblance is not descent. \emph{Río} and \emph{river} have four letters in common and nothing else.
\end{cousinweb}

\begin{grammarlens}[title={What you've built}]
You can now name the water in the landscape: \emph{el río está en el valle}.
\end{grammarlens}

\subsection*{Your turn}
Say these aloud:

\begin{itemize}
\raggedright
  \item \emph{el río}
  \item \emph{un río}
  \item \emph{el río está en el valle}
\end{itemize}

\begin{tcolorbox}[breakable,colback=teal!4,colframe=teal!35!black,title={Before you move on}]
What is a rival, taken literally? (Someone drawing on the same stream.) And is English \emph{river} the same word as \emph{río}? (No — it comes from \emph{rīpa}, a bank, like \emph{arrive}.)
\end{tcolorbox}
\section[la arena]{la arena — the sand that gave its name to the building}

\label{lesson:ES-C369-arena}



\begin{quote}
Say \emph{the bull}, then \emph{the river}. (\emph{El toro}; \emph{el río}.)
\end{quote}

\subsection*{What to know first}
\begin{itemize}
\raggedright
  \item el mar — what puts it there.
  \item la orilla — the strip it covers.
\end{itemize}

\begin{sounds}
\begin{itemize}
\raggedright
  \item \emph{a-RE-na}, three beats with the weight on the second. The \emph{r} between vowels is a single tap. $\to$ pronunciation reference
\end{itemize}
\end{sounds}

\begin{cousinweb}
\textbf{la arena} — \textquotedblleft{}sand.\textquotedblright{}

Latin \textbf{harēna}, sand — and after that the trail goes cold. This is one of the words where the reference books shrug.

The standing guess is Etruscan, the language of the people Rome grew up beside, on the grounds that the ending looks Etruscan. The other piece of evidence is a form \textbf{fasēna}, which the Roman scholar Varro reports as the Sabine version of the word. Neither has ever been made to prove anything, and the verdict printed in the best modern dictionary of Latin origins is: unknown.

So keep the word and let the ancestry go. It is worth meeting a word that nobody can trace, because most vocabularies contain a few and it is honest to say which ones they are.

What the word did next is well documented. A Roman amphitheatre had \textbf{sand} spread across its floor, to give footing and to soak up what got spilled on it. People called the floor by the name of what was on it. Then they called the whole building that. English \textbf{arena} is that Latin word for sand, and every political arena, every sporting arena, is named after a surface.

Spanish kept both jobs. \emph{La arena} is the beach, and \emph{la arena} is the bullring.
\end{cousinweb}

\begin{grammarlens}[title={What you've built}]
You can now name what the beach is made of: \emph{la arena}.
\end{grammarlens}

\subsection*{Your turn}
Say these aloud:

\begin{itemize}
\raggedright
  \item \emph{la arena}
  \item \emph{la arena del mar}
  \item \emph{hay arena en la orilla}
\end{itemize}

\begin{tcolorbox}[breakable,colback=teal!4,colframe=teal!35!black,title={Before you move on}]
What was spread on the floor of a Roman amphitheatre? (Sand.) And what is the honest answer about where \emph{harēna} came from? (Nobody knows.)
\end{tcolorbox}
\section[el polvo]{el polvo — one Latin word, two English roads}

\label{lesson:ES-C369-polvo}



\begin{quote}
Say \emph{the river}, then \emph{sand}. (\emph{El río}; \emph{la arena}.)
\end{quote}

\subsection*{What to know first}
\begin{itemize}
\raggedright
  \item la tierra — what it comes off.
  \item limpio — its opposite, in one word.
\end{itemize}

\begin{sounds}
\begin{itemize}
\raggedright
  \item \emph{POL-vo}, two beats with the weight on the first. The \emph{v} is pronounced the same as a Spanish \emph{b}, with the lips barely meeting. $\to$ pronunciation reference
\end{itemize}
\end{sounds}

\begin{cousinweb}
\textbf{el polvo} — \textquotedblleft{}dust.\textquotedblright{}

Latin \textbf{pulvis}, with the stem \emph{pulver-}: dust, powder, the fine stuff left when something has been ground or has fallen apart.

English took that word twice, and the two arrivals do not look alike:

\begin{itemize}
\raggedright
  \item \textbf{powder} came through Old French \emph{poudre}. French had put a \textbf{d} into it that Latin never had — the sort of consonant that grows between an \emph{l} and an \emph{r} because it is easier to say than to avoid;
  \item \textbf{pulverize} was taken off the written Latin stem centuries later, so it still has the \emph{v} and no \emph{d} at all.
\end{itemize}

German kept both and gave each a separate job, which makes the split easy to hold on to: \textbf{Puder} is the powder for a face, borrowed from French with the \emph{d}; \textbf{Pulver} is gunpowder, taken from the Latin with the \emph{v}. Same word, two doors, two meanings, no confusion.

One connection you may be offered elsewhere is worth declining. \textbf{Pollen} is sometimes filed with \emph{pulvis}, on the grounds that both are fine ground stuff. Specialists doubt it, and the objection is a good one: flour and dust are what you get at the two ends of the same milling, and words for them do not reliably share ancestors.
\end{cousinweb}

\begin{grammarlens}[title={What you've built}]
You can now describe what settles on a shelf: \emph{hay polvo}.
\end{grammarlens}

\subsection*{Your turn}
Say these aloud:

\begin{itemize}
\raggedright
  \item \emph{el polvo}
  \item \emph{hay polvo}
  \item \emph{la mesa tiene polvo}
\end{itemize}

\begin{tcolorbox}[breakable,colback=teal!4,colframe=teal!35!black,title={Before you move on}]
Which of \emph{powder} and \emph{pulverize} came through French? (\emph{Powder} — and French added the \emph{d}.) And which German word is gunpowder? (\emph{Pulver}.)
\end{tcolorbox}
\section[el humo]{el humo — an f that turned into a silent h}

\label{lesson:ES-C369-humo}



\begin{quote}
Say \emph{sand}, then \emph{dust}. (\emph{La arena}; \emph{el polvo}.)
\end{quote}

\subsection*{What to know first}
\begin{itemize}
\raggedright
  \item el fuego — what makes it.
  \item la nube — what it looks like from far off.
\end{itemize}

\begin{sounds}
\begin{itemize}
\raggedright
  \item \emph{HU-mo}, two beats with the weight on the first, and the \emph{h} is completely silent — the word begins on the \emph{u}. $\to$ pronunciation reference
\end{itemize}
\end{sounds}

\begin{cousinweb}
\textbf{el humo} — \textquotedblleft{}smoke.\textquotedblright{}

Latin \textbf{fūmus}, smoke. Spanish writes \textbf{humo}, with a silent \emph{h} where the \emph{f} was.

That \emph{h} is a gravestone. At one stage Spanish said \emph{fumo}, then softened the \emph{f} to a breath, then stopped pronouncing the breath, and kept writing the letter anyway. Medieval texts catch the word mid-fall: \emph{fumo} on the page long after speakers had begun to let it go. The same change carried Latin \emph{facere} to \textbf{hacer} and \emph{filium} to \textbf{hijo}.

English never had that change, so English still has the \emph{f} wherever it took the word off the Latin:

\begin{itemize}
\raggedright
  \item \textbf{fume} — the vapour itself, and by extension the temper of somebody fuming;
  \item \textbf{fumigate} — to treat a place with smoke;
  \item \textbf{fumitory} — a small plant whose name means \emph{fūmus terrae}, smoke of the earth, for its grey-green haze low to the ground.
\end{itemize}

The best of them is \textbf{perfume}: \emph{per} plus the smoke, \textquotedblleft{}through smoke.\textquotedblright{} A perfume was what rose off something burning, and the first English uses of the word mean precisely that — the fumes of a burning substance, incense in a room, not a bottle on a shelf. The bottle came later and kept the word for smoke.
\end{cousinweb}

\begin{grammarlens}[title={What you've built}]
You can now say what a fire sends up: \emph{hay humo}.
\end{grammarlens}

\subsection*{Your turn}
Say these aloud:

\begin{itemize}
\raggedright
  \item \emph{el humo}
  \item \emph{hay humo}
  \item \emph{el fuego hace humo}
\end{itemize}

\begin{tcolorbox}[breakable,colback=teal!4,colframe=teal!35!black,title={Before you move on}]
Which Latin letter is the \emph{h} of \emph{humo} standing in for? (\emph{F}.) And what did a perfume originally mean? (Something that rose off a fire — through smoke.)
\end{tcolorbox}
\section[la ceniza]{la ceniza — and the cinder that is wearing a costume}

\label{lesson:ES-C369-ceniza}



\begin{quote}
Say \emph{dust}, then \emph{smoke}. (\emph{El polvo}; \emph{el humo}.)
\end{quote}

\subsection*{What to know first}
\begin{itemize}
\raggedright
  \item el fuego — what leaves it behind.
  \item el humo — the other thing a fire sends out.
\end{itemize}

\begin{sounds}
\begin{itemize}
\raggedright
  \item \emph{ce-NI-za}, three beats with the weight on the second. Both the \emph{c} and the \emph{z} are the same sound in one accent or another — a soft \emph{th} in much of Spain, an \emph{s} across the Americas. $\to$ pronunciation reference
\end{itemize}
\end{sounds}

\begin{cousinweb}
\textbf{la ceniza} — \textquotedblleft{}ash.\textquotedblright{}

Latin \textbf{cinis}, with the stem \emph{ciner-}: ashes. A late spoken form built on it gave Spanish \emph{ceniza}. English has the same Latin word wherever it borrowed off the page — to \textbf{incinerate} is to reduce to ashes, and a \textbf{cinerary} urn holds what is left after one.

So English \textbf{cinder} must belong here too. It does not.

\textbf{Cinder} is Old English \emph{sinder}, and it meant the \textbf{slag} that comes off iron — a Germanic word, no Latin in it anywhere. It began with an \emph{s}. Then French \emph{cendre}, the genuine descendant of \emph{cinis}, was sitting nearby in the same language, meaning nearly the same thing, and English respelled its own native word with a \emph{c} to match. The costume has been on for so long that the word looks Latin to everybody who meets it.

That is worth holding on to, because it is the reverse of the usual mistake. Words that look unrelated often are related. This is one that looks related and is not, and the resemblance was applied on purpose by scribes.

The girl in the ashes carries the honest version. French calls her \textbf{Cendrillon}, from \emph{cendre}; Spanish calls her \textbf{Cenicienta}, from \emph{ceniza}; German calls her \textbf{Aschenputtel}, from \emph{Asche}. Three languages naming her after the hearth she sleeps by, each out of its own word for ash.
\end{cousinweb}

\begin{grammarlens}[title={What you've built}]
You can now name what a fire leaves: \emph{la ceniza}.
\end{grammarlens}

\subsection*{Your turn}
Say these aloud:

\begin{itemize}
\raggedright
  \item \emph{la ceniza}
  \item \emph{la ceniza del fuego}
  \item \emph{hay ceniza}
\end{itemize}

\begin{tcolorbox}[breakable,colback=teal!4,colframe=teal!35!black,title={Before you move on}]
Is English \emph{cinder} from the same word as \emph{ceniza}? (No — it is Germanic, only respelled to look Latin.) And what is \emph{Cenicienta} named after? (\emph{La ceniza}, the ash.)
\end{tcolorbox}
